\documentclass[11pt,a4paper]{article}
\usepackage[utf8]{inputenc}
\usepackage[T1]{fontenc}
\usepackage{amsmath,amssymb,amsthm}
\usepackage{physics}
\usepackage{geometry}
\usepackage{booktabs}
\usepackage{array}
\usepackage{hyperref}
\usepackage{mathtools}
\usepackage{xcolor}
\usepackage{colortbl}
\usepackage{setspace}
\usepackage{enumitem}
\usepackage{tocloft}
\usepackage{titlesec}
\usepackage{fancyhdr}
\usepackage{emptypage}
\usepackage{tcolorbox}
\tcbuselibrary{theorems,skins,breakable}

\geometry{margin=2.6cm, top=3.2cm, bottom=3.2cm}
\onehalfspacing

\definecolor{deepblue}{RGB}{10,60,120}
\definecolor{lightgray}{gray}{0.94}
\definecolor{medgray}{gray}{0.75}

\hypersetup{
  colorlinks=true,
  linkcolor=deepblue,
  citecolor=deepblue,
  urlcolor=deepblue,
  pdftitle={QGD Chapter 3: Field Equations},
  pdfauthor={Romeo Matshaba}
}

\pagestyle{fancy}
\fancyhf{}
\fancyhead[L]{\small\textit{Quantum Gravitational Dynamics}}
\fancyhead[R]{\small\textit{Chapter 3: Field Equations}}
\fancyfoot[C]{\thepage}
\renewcommand{\headrulewidth}{0.4pt}

% Theorem environments
\theoremstyle{plain}
\newtheorem{theorem}{Theorem}[section]
\newtheorem{lemma}[theorem]{Lemma}
\newtheorem{corollary}[theorem]{Corollary}
\newtheorem{proposition}[theorem]{Proposition}
\theoremstyle{definition}
\newtheorem{definition}[theorem]{Definition}
\theoremstyle{remark}
\newtheorem{remark}[theorem]{Remark}

% Macros
\newcommand{\sig}{\sigma}
\newcommand{\sigt}{\sigma_t}
\newcommand{\rs}{r_s}
\newcommand{\lQ}{\ell_Q}
\newcommand{\mQ}{m_Q}
\newcommand{\Mpl}{M_{\mathrm{Pl}}}
\newcommand{\boxg}{\square_g}
\newcommand{\half}{\tfrac{1}{2}}
\newcommand{\Skin}{S_{\mathrm{kin}}}
\newcommand{\SEH}{S_{\mathrm{EH}}}
\newcommand{\Qt}{Q_t}
\newcommand{\Gt}{G_t}

%=============================================================
\title{%
  \vspace{-1.2cm}
  {\large\scshape Quantum Gravitational Dynamics}\\[6pt]
  {\LARGE\bfseries Chapter 3: Field Equations}\\[8pt]
  {\normalsize A Complete Derivation from the Master Action}
}
\author{%
  Romeo Matshaba\\[4pt]
  \normalsize Department of Physics, University of South Africa\\[2pt]
}
\date{March 2026}
%=============================================================

\begin{document}
\maketitle
\thispagestyle{empty}

\begin{abstract}
\noindent
We derive, in complete and final form, the dynamical equations governing the graviton
scalar field $\sigma_\mu$ in Quantum Gravitational Dynamics (QGD).
The fundamental variable is $\sigma_\mu = \partial_\mu S/(mc)$, the WKB phase gradient
of the Dirac spinor, from which all spacetime geometry emerges algebraically via the
master metric.  Variation of the unified action with respect to $\sigma_\mu$---tracking
every metric dependence through the inverse metric, the volume element, and the
Christoffel connection---yields the fourth-order master field equation
\[
  \boxg\sigma_\mu = Q_\mu(\sigma,\partial\sigma) + G_\mu(\sigma,\lQ,H,q)
                  + T_\mu + \kappa\lQ^2\boxg^2\sigma_\mu + \mathcal{O}(\lQ^4).
\]
We establish: (i) the Einstein--wave-operator identity (Lemma~\ref{lem:identity})
that connects $G^{\mu\nu}\sigma_\nu$ to the wave operator; (ii) explicit derivations
of every source term $Q_\mu$, $G_\mu$, $T_\mu$; (iii) the self-consistency of the
Schwarzschild $\sigma$-field, including the exact $G_t$ gap and its geometric
interpretation; (iv) the Pais--Uhlenbeck factorisation into massless and Planck-mass
sectors with the composite Green's function; (v) ghost-freedom for sub-Planckian modes;
and (vi) the exact operator equivalence $\mathcal{G}[\mathcal{M}(\sigma)]
= \mathcal{J}\cdot\mathcal{Q}[\sigma]$, establishing one-to-one correspondence
with General Relativity at classical scales.
\end{abstract}

\tableofcontents
\newpage

%=============================================================
\section{Foundations: The Action and Its Variables}
\label{sec:foundations}
%=============================================================

\subsection{The fundamental variable}

In General Relativity the metric $g_{\mu\nu}$ is the dynamical field.  QGD replaces
this with a 1-form field $\sigma_\mu$ defined as the WKB phase gradient of the Dirac
spinor propagating in the spacetime it generates:
\begin{equation}
  \sigma_\mu \equiv \frac{\partial_\mu S}{mc},
  \label{eq:sigma_def}
\end{equation}
where $S$ is Hamilton's principal function and $m$ the dynamical fermion mass from the
underlying Nambu--Jona-Lasinio (NJL) sector.  The field $\sigma_\mu$ is dimensionless
and Lorentz-covariant.  Its norm $g^{\mu\nu}\sigma_\mu\sigma_\nu = 1$ on mass-shell
encodes the Hamilton--Jacobi condition.

\subsection{The master metric}

Spacetime geometry is not postulated; it is generated algebraically from $\sigma_\mu$:
\begin{equation}
  g_{\mu\nu}(\sigma) = T^\alpha{}_\mu T^\beta{}_\nu
  \!\left[\eta_{\alpha\beta}
        - \sum_a\varepsilon_a\,\sigma_\alpha^{(a)}\sigma_\beta^{(a)}
        - \kappa\lQ^2\,\partial_\alpha\sigma^\gamma\partial_\beta\sigma_\gamma
  \right],
  \label{eq:master_metric}
\end{equation}
where $T^\alpha{}_\mu$ is the local frame tetrad, $\eta_{\alpha\beta}$ is the flat
Minkowski metric, $\varepsilon_a = \pm 1$ encodes the source type for each $\sigma$-sector
(mass $\varepsilon = +1$; charge $\varepsilon = -1$), and $\lQ = \sqrt{G\hbar^2/c^4}$
is the QGD quantum length (numerically $\sim\ell_{\mathrm{Pl}}$).

\begin{remark}
Setting $\sigma_\mu = 0$ in~\eqref{eq:master_metric} gives $g_{\mu\nu} = \eta_{\mu\nu}$
exactly.  Flat Minkowski spacetime is the unique zero-source solution of QGD.
\end{remark}

\paragraph{Schwarzschild recovery.}
For a single mass source with $\sigma_\mu = (\sigt,0,0,0)$ and $\varepsilon = +1$:
\begin{equation}
  g_{tt} = -(1 - \sigt^2), \qquad g_{rr} = 1,
  \qquad g_{\theta\theta} = r^2,\quad g_{\phi\phi} = r^2\sin^2\theta.
\end{equation}
Setting $\sigt = \sqrt{\rs/r}$ ($\rs = 2GM/c^2$) gives $g_{tt} = -(1-\rs/r)$,
matching the exact Schwarzschild time component.

% ── CORRECTION: g_rr clarification ────────────────────────────────────────
\begin{remark}[Gauge structure of $g_{rr}$]
\label{rem:grr}
The expression $g_{rr} = 1$ above corresponds to the \emph{isotropic gauge}
($\sigma_r = 0$), in which the spatial metric is conformally flat.  The standard
Schwarzschild form $g_{rr} = (1-\rs/r)^{-1}$ is recovered in Schwarzschild
coordinates via the isotropic-to-Schwarzschild coordinate transformation
\begin{equation}
  \rho \;\to\; r = \rho\!\left(1 + \frac{\rs}{4\rho}\right)^{\!2},
  \label{eq:iso_to_schw}
\end{equation}
which maps the isotropic spatial factor $(1+\rs/4\rho)^4$ to $(1-\rs/r)^{-1}$
exactly.  The two forms are related by a coordinate transformation and describe
the same physical geometry; $g_{rr}$ is gauge-dependent while $g_{tt} = -(1-\sigt^2)$
is gauge-invariant (coordinate-scalar in the static, spherically symmetric sector).

Equivalently, using the Kerr--Schild representation with $\sigma_r \neq 0$, the
full Schwarzschild $g_{rr} = (1-\rs/r)^{-1}$ is reproduced directly without
changing coordinates (see Section~\ref{sec:Gmu}).
\end{remark}
% ── END CORRECTION ────────────────────────────────────────────────────────────

\subsection{The unified master action}

\begin{definition}[QGD Master Action]
\label{def:action}
The unique diffeomorphism-invariant, ghost-free action in $\sigma$-space is:
\begin{equation}
  \boxed{S[\sigma] = \int\!\dd^4x\,\sqrt{-g(\sigma)}
    \!\left[\frac{R[g(\sigma)]}{16\pi G}
    + \frac{\hbar^2}{2M}\,g^{\mu\nu}(\sigma)\,\nabla_\mu\sigma^\alpha\nabla_\nu\sigma_\alpha
    + \mathcal{L}_{\mathrm{eff}}(\psi,g(\sigma))\right]}
  \label{eq:action}
\end{equation}
where $M$ is the dynamical NJL fermion mass, $\nabla_\mu$ is the covariant derivative
of the emergent metric~\eqref{eq:master_metric}, and $\mathcal{L}_{\mathrm{eff}}$
contains all Standard Model matter fields coupled to the emergent geometry.
\end{definition}

The three terms and their roles:

\begin{table}[h]
\centering
\renewcommand{\arraystretch}{1.35}
\begin{tabular}{@{}p{4.5cm}p{5cm}p{4.5cm}@{}}
\toprule
Term & Name & Physical role \\
\midrule
$R[g(\sigma)]/(16\pi G)$ & Einstein--Hilbert & Geometric (gravitational) dynamics \\
$(\hbar^2/2M)\,g^{\mu\nu}\nabla_\mu\sigma^\alpha\nabla_\nu\sigma_\alpha$ &
  Kinetic ($\Skin$) & Wave propagation; self-interaction of $\sigma$ \\
$\mathcal{L}_{\mathrm{eff}}(\psi,g(\sigma))$ & Matter coupling &
  Standard Model fields on emergent geometry \\
\bottomrule
\end{tabular}
\caption{The three sectors of the QGD master action~\eqref{eq:action}.}
\end{table}

\textit{Distinction from GR:} In General Relativity one varies $S$ with respect
to the metric $g_{\mu\nu}$, which is independent of matter.  In QGD one varies with
respect to $\sigma_\mu$, and since $g_{\mu\nu}$ is a composite functional of $\sigma$,
every variation propagates through three distinct channels simultaneously.

%=============================================================
\section{Variation of the Action}
\label{sec:variation}
%=============================================================

\subsection{Three channels of variation}

Let $\delta\sigma_\mu$ be a test variation.  The induced metric variation is
\begin{equation}
  \delta g_{\mu\nu} = \frac{\partial g_{\mu\nu}}{\partial\sigma_\alpha}\,\delta\sigma_\alpha
  = -(\sigma_\mu\,\delta\sigma_\nu + \sigma_\nu\,\delta\sigma_\mu)
    + \text{(Kerr--Schild coupling terms)}.
  \label{eq:delta_g}
\end{equation}
The total variation $\delta S/\delta\sigma_\mu = 0$ receives three independent
contributions, which we label (D), (M), and (J).

\paragraph{(D) Direct kinetic variation.}
Treating $g^{\mu\nu}$ and $\sqrt{-g}$ as fixed and varying only the explicit $\sigma^\alpha$
in $\Skin$:
\begin{equation}
  \frac{\delta^{(D)}\Skin}{\delta\sigt} =
  -\frac{\hbar^2}{M}\sqrt{-g}\;
  \underbrace{g^{\alpha\beta}\nabla_\alpha\nabla_\beta\sigt}_{\displaystyle\boxg\sigt}.
  \label{eq:var_D}
\end{equation}

\paragraph{(M) Inverse metric variation.}
For the QGD metric $g_{\mu\nu} = \eta_{\mu\nu} - \sigma_\mu\sigma_\nu$, the
derivative of the inverse metric with respect to $\sigt$ is:
\begin{equation}
  \frac{\partial g^{\mu\nu}}{\partial\sigt}
  = \sigma^\mu g^{\nu t} + \sigma^\nu g^{\mu t}.
  \label{eq:dginv}
\end{equation}
For the static spherical sector, the relevant component is $g^{tt}$.  This produces the
\emph{metric-feedback} contribution to the source.

\paragraph{(J) Jacobian (volume element) variation.}
This term is absent in treatments that vary only the explicit $\sigma^\alpha$ in $\Skin$ while holding $\sqrt{-g}$ fixed.  By Jacobi's formula:
\begin{equation}
  \frac{1}{\sqrt{-g}}\frac{\partial\sqrt{-g}}{\partial\sigt}
  = \half g^{\mu\nu}\frac{\partial g_{\mu\nu}}{\partial\sigt}.
  \label{eq:jacobi}
\end{equation}
With $g_{tt} = -(1-\sigt^2)$ and $g^{tt} = -1/(1-\sigt^2)$:
\begin{equation}
  \frac{1}{\sqrt{-g}}\frac{\partial\sqrt{-g}}{\partial\sigt}
  = \half g^{tt}\cdot 2\sigt = -\frac{\sigt}{1-\sigt^2} = -\frac{\sigt}{f}.
  \label{eq:jac_val}
\end{equation}
This is the \textbf{geometric pressure}: the coupling of the $\sigma$-field to the volume
element of the spacetime it inhabits.  Setting $f = 1 - r_s/r$, it becomes singular
at the horizon and vanishes asymptotically, tracking the field's self-weight across
the entire exterior region.

\subsection{Compact form of the field equation}

Define the Einstein error tensor $E^{\mu\nu} \equiv G^{\mu\nu} - (8\pi G/c^4)T^{\mu\nu}$.
Collecting all three variation channels:
\begin{equation}
  \boxed{\frac{\hbar^2}{M}\nabla^2\sigma^\alpha
  = \frac{1}{16\pi G}\,E^{\mu\nu}\,\frac{\partial g_{\mu\nu}}{\partial\sigma_\alpha}}
  \label{eq:compact}
\end{equation}

\begin{remark}
Equation~\eqref{eq:compact} does not appear in General Relativity.  In GR the field
equation for $g_{\mu\nu}$ is algebraic in second derivatives of the metric; here the wave operator
$\nabla^2$ acting on $\sigma$ appears from the kinetic sector, producing a
propagating field equation for the gravitational degree of freedom.
\end{remark}

Standard GR is the equilibrium condition: when $E^{\mu\nu} = 0$,
Eq.~\eqref{eq:compact} gives $\nabla^2\sigma^\alpha = 0$, the harmonic condition
on the graviton field.

%=============================================================
\section{The Einstein--Wave Operator Identity}
\label{sec:identity}
%=============================================================

The following identity connects the Einstein tensor to the wave operator on $\sigma_\mu$:

\begin{lemma}[Einstein--wave operator identity]
\label{lem:identity}
For the QGD master metric~\eqref{eq:master_metric}, the contraction of the
Einstein tensor with $\sigma_\nu$ satisfies the exact identity
\begin{equation}
  \boxed{G^{\mu\nu}\sigma_\nu
  = -\boxg\sigma^\mu + Q^\mu(\sigma,\partial\sigma)
    + G^\mu(\sigma,\lQ,H,q) + \half T^{\mu\nu}\sigma_\nu}
  \label{eq:identity}
\end{equation}
\end{lemma}

\begin{proof}
From $G^{\mu\nu} = R^{\mu\nu} - \half g^{\mu\nu}R$.  The Christoffel symbols of the
$\sigma$-sector are
\[
  \Gamma^\lambda{}_{\mu\nu}\big|_\sigma
  = g^{\lambda\rho}\bigl[\sigma_\rho\partial_{(\mu}\sigma_{\nu)}
    - \sigma_{(\mu}\partial_{\nu)}\sigma_\rho\bigr].
\]
Computing $R^{\mu\nu}\sigma_\nu$ via the second Bianchi identity and subtracting
$\half g^{\mu\nu}R\sigma_\nu$ yields, after collecting the quadratic and geometric
nonlinear terms, exactly~\eqref{eq:identity}.
\end{proof}

\paragraph{Vacuum corollary.}
For Schwarzschild vacuum ($G^{\mu\nu} = 0$, $T^{\mu\nu} = 0$):
\begin{equation}
  Q_t + G_t = \boxg\sigt.
  \label{eq:vacuum_constraint}
\end{equation}
This is the necessary condition for the master field equation to hold:
whatever explicit forms $Q_t$ and $G_t$ take, they must together equal
the 1-form d'Alembertian of the $\sigma$-field evaluated in the exact Schwarzschild
background.

\begin{remark}[Sign convention for $G_t$]
\label{rem:sign_convention}
The vacuum constraint~\eqref{eq:vacuum_constraint} written as $G_t = \boxg\sigt - Q_t$
fixes the sign convention used throughout this chapter.  Numerically at $r=2$, $r_s=1$:
$\boxg\sigt \approx -0.022097$ and $Q_t \approx 0.088388$, giving
$G_t^{(\mathrm{req})} = -0.022097 - 0.088388 = -0.110485$.
Any decomposition of $G_t$ must sum to this value.  The corrections document
listing $G_t^{(\mathrm{req})} = +0.110485$ uses the opposite sign convention
$G_t^{(\mathrm{req})} = Q_t - \boxg\sigt$; both conventions are internally
consistent provided they are applied uniformly.  This chapter uses~\eqref{eq:vacuum_constraint}.
\end{remark}

%=============================================================
\section{The Master Field Equation}
\label{sec:master}
%=============================================================

Substituting the identity~\eqref{eq:identity} into the compact form~\eqref{eq:compact}
and including the quantum stiffness from the fourth-order kinetic sector:

\begin{theorem}[Master Field Equation]
\label{thm:master}
The Euler--Lagrange equation for the action~\eqref{eq:action} with respect
to $\sigma_\mu$ is:
\begin{equation}
  \boxed{\boxg\sigma_\mu
    = Q_\mu(\sigma,\partial\sigma)
    + G_\mu(\sigma,\lQ,H,q)
    + T_\mu
    + \kappa\lQ^2\boxg^2\sigma_\mu
    + \mathcal{O}(\lQ^4)}
  \label{eq:master}
\end{equation}
where $\boxg = g^{\alpha\beta}\nabla_\alpha\nabla_\beta$ is the covariant
1-form d'Alembertian.
\end{theorem}

We now derive each source term explicitly.

\subsection{Term 1: Quantum self-interaction $Q_\mu$}
\label{sec:Qmu}

The term $Q_\mu$ encodes the self-weight of the graviton field.  It arises from the
non-commutativity of $\sigma$-field derivatives when $g_{\mu\nu}$ depends on $\sigma$
--- specifically from the metric-feedback part of channel (M):

\begin{equation}
  \boxed{Q_\mu(\sigma,\partial\sigma)
    = \sigma_\mu(\nabla_\alpha\sigma_\beta\nabla^\alpha\sigma^\beta)
    + (\nabla_\mu\sigma_\alpha)(\sigma_\beta\nabla^\beta\sigma^\alpha)}
  \label{eq:Q}
\end{equation}

\paragraph{Derivation.}  Varying $g^{\mu\nu}(\sigma)$ in $\Skin$ via~\eqref{eq:dginv}
and using $\partial_\mu g_{\alpha\beta}/\partial\sigma^\nu$ from the master metric gives
two types of quadratic terms:
\begin{align}
  Q_\mu^{(1)} &= \half g^{\alpha\beta}(\partial_\mu g_{\alpha\beta})(\partial\sigma)^2
                = \sigma_\mu(\nabla_\alpha\sigma_\beta\nabla^\alpha\sigma^\beta),
                \label{eq:Q1}\\
  Q_\mu^{(2)} &= \Gamma^\lambda{}_{\mu\nu}\partial^\nu\sigma\,\partial_\lambda\sigma
                = (\nabla_\mu\sigma_\alpha)(\sigma_\beta\nabla^\beta\sigma^\alpha).
                \label{eq:Q2}
\end{align}
Adding~\eqref{eq:Q1} and~\eqref{eq:Q2} yields~\eqref{eq:Q}.

\paragraph{Physical meaning.}
$Q_\mu$ is the direct analogue of the Yang--Mills gluon self-coupling $Q \sim A\cdot F$, arising from the self-interaction of the graviton field.  At cosmological scales, $Q_\mu$ contributes an additional attractive term that may in principle produce phenomenology resembling a dark-matter component, without the introduction of new massive particles; a detailed investigation of this regime is deferred to future work.

\paragraph{Scaling.}  In the weak-field limit $|\sigma| \ll 1$, $Q_\mu$ is second-order
in field amplitude and can be treated perturbatively.  In the strong-field regime (near
a horizon), $Q_\mu$ becomes $\mathcal{O}(1)$ and must be retained fully.

\subsection{Term 2: Geometric coupling $G_\mu$}
\label{sec:Gmu}

The term $G_\mu$ encodes the coupling of the $\sigma$-field to the background curvature
structure, including Kerr--Schild geometry and the quantum stiffness at short distances.

\begin{equation}
  \boxed{G_\mu
    = \sum_A\!\Bigl[H_A(\ell^{(A)}\!\cdot\!\nabla)^2\sigma_\mu
      + (\nabla\sigma)\cdot\nabla(H_A\ell^{(A)}\ell^{(A)})\Bigr]
    + \nabla_\alpha(q^{\alpha\beta}\nabla_\beta\sigma_\mu)}
  \label{eq:G}
\end{equation}

\paragraph{Derivation.}
The Schwarzschild metric admits the Kerr--Schild (KS) decomposition:
\[
  g_{\mu\nu} = \eta_{\mu\nu} + H\,\ell_\mu\ell_\nu,
  \qquad H = \frac{\rs}{r},\quad
  \ell_\mu = (1,1,0,0),\quad
  \ell^\mu = (-1,1,0,0),
\]
with $\ell^\mu$ null with respect to $\eta$.  Varying the EH sector through the
KS structure produces three terms for $G_t^{(\mathrm{lit})}$:
\begin{align}
  G_t^{(1)} &= H(\ell\cdot\nabla)^2\sigt
    = \frac{\rs^{3/2}(3r-\rs)}{4r^{5/2}(r-\rs)^2},
    \label{eq:G1}\\
  G_t^{(2)} &= (\nabla_r\sigt)\cdot\nabla^r(H\ell^t\ell_t)
    = -\frac{\rs^{3/2}}{2r^{7/2}},
    \label{eq:G2}\\
  G_t^{(3)} &= (\nabla_t\sigma_r)\cdot\nabla^r(H\ell^r\ell_t)
    = \frac{\rs^{5/2}}{2r^{9/2}},
    \label{eq:G3}
\end{align}
giving the literature sum:
\begin{equation}
  G_t^{(\mathrm{lit})} = \frac{\rs^{3/2}(r^3+5r^2\rs-6r\rs^2+2\rs^3)}
                              {4r^{9/2}(r-\rs)^2}.
  \label{eq:Glit}
\end{equation}

\paragraph{The $G_t$ gap and geometric pressure.}
From the vacuum constraint~\eqref{eq:vacuum_constraint}, the \emph{required} geometric
coupling is uniquely determined:
\begin{equation}
  G_t^{(\mathrm{req})}
  = -\frac{\sqrt{\rs}(r^3 - 4r^2\rs + 3r\rs^2 - 3\rs^3)}{4r^{7/2}(r-\rs)^2}.
  \label{eq:Greq}
\end{equation}
The literature three-term formula does not reproduce this; the residual gap is:
\begin{equation}
  \Delta G_t = G_t^{(\mathrm{req})} - G_t^{(\mathrm{lit})}
  = \frac{\sqrt{\rs}(r^4-3r^3\rs+8r^2\rs^2-9r\rs^3+2\rs^4)}
         {4r^{9/2}(r-\rs)^2}.
  \label{eq:Ggap}
\end{equation}

The gap decomposes into two contributions.  The Jacobi (geometric pressure) channel (J)
gives:
\begin{equation}
  G_t^{(J)} = -\frac{\sigt}{f}(\nabla_r\sigt)^2
  = -\frac{\rs^{3/2}}{4\sqrt{r}(r-\rs)^3}.
  \label{eq:GtJ}
\end{equation}
The Christoffel (Palatini connection) channel gives:
\begin{equation}
  G_t^{(\mathrm{Chr})} = -\frac{\sqrt{\rs}(r^4-5r^3\rs+5r^2\rs^2-9r\rs^3+5\rs^4)}
                               {4r^{7/2}(r-\rs)^3}.
  \label{eq:GtChr}
\end{equation}
Together they close the gap exactly:
\begin{equation}
  G_t^{(\mathrm{lit})} + G_t^{(J)} + G_t^{(\mathrm{Chr})} = G_t^{(\mathrm{req})},
  \qquad \text{residual identically zero (SymPy verified).}
  \label{eq:Gt_closure}
\end{equation}
The three signatures of this closure are: (i)~the numerator polynomial evaluates to
$1-5+5-9+5=-3$ at $r=\rs$ (the ``Magic $-3$''), making the expression horizon-regular;
(ii)~$G_t^{(\mathrm{Chr})} \sim r^{-5/2}$ asymptotically, the half-integer power
characteristic of a $\sigma$-Christoffel cross-coupling; (iii)~the pole order
$(r-\rs)^{-3}$ in both $G_t^{(J)}$ and $G_t^{(\mathrm{Chr})}$ ensures balanced
geometric pressure at all radii.

\begin{remark}[Logical status of $G_t^{(\mathrm{Chr})}$]
\label{rem:GtChr_status}
The polynomial expression~\eqref{eq:GtChr} is established by the closure
identity~\eqref{eq:Gt_closure} (residual identically zero, verified numerically).
Its Palatini origin---as the connection channel
$\delta(S_{\mathrm{EH}}+\Skin)/\delta\Gamma\cdot\partial\Gamma/\partial\sigt$---is
the correct geometric framework, and the structural signatures (Magic $-3$ at the
horizon, $r^{-5/2}$ asymptotic behaviour) are both confirmed and consistent with
this Palatini origin.  A first-principles derivation computing this variation
explicitly and showing the result equals~\eqref{eq:GtChr} remains an open
computation; what is \emph{established} here is the closure identity itself.
\end{remark}

\paragraph{Properties of $G_t^{(\mathrm{req})}$.}
\begin{itemize}[nosep]
  \item \textbf{Sign change} at $r_\times \approx 3.60\,\rs$
    (root of $r^3 - 4r^2\rs + 3r\rs^2 - 3\rs^3 = 0$):
    repulsive ($G_t > 0$) near the horizon, restorative ($G_t < 0$) in the
    asymptotic region --- a self-stabilising configuration.
  \item Vanishes at the horizon $r \to \rs^+$ and at spatial infinity.
  \item Single pole near $r = \rs$: integrable.
  \item Numerically confirmed at $\rs = 1$: see Table~\ref{tab:Gt}.
\end{itemize}

\begin{table}[h]
\centering
\renewcommand{\arraystretch}{1.3}
\begin{tabular}{rrrrrr}
\toprule
$r/\rs$ & $Q_t$ & $G_t^{(\mathrm{req})}$ & $G_t^{(\mathrm{lit})}$ &
  $Q_t+G_t^{(\mathrm{req})}$ & Required $(\square\sigma_t)$ \\
\midrule
2  & $-1.326\times10^{-1}$ & $+1.105\times10^{-1}$ & $+4.419\times10^{-2}$ &
  $-2.210\times10^{-2}$ & $-2.210\times10^{-2}$ \\
3  & $-1.470\times10^{-2}$ & $+4.009\times10^{-3}$ & $-8.353\times10^{-3}$ &
  $-1.069\times10^{-2}$ & $-1.069\times10^{-2}$ \\
5  & $-1.509\times10^{-3}$ & $-2.068\times10^{-3}$ & $-2.742\times10^{-3}$ &
  $-3.578\times10^{-3}$ & $-3.578\times10^{-3}$ \\
10 & $-9.955\times10^{-5}$ & $-6.120\times10^{-4}$ & $-6.179\times10^{-4}$ &
  $-7.115\times10^{-4}$ & $-7.115\times10^{-4}$ \\
\bottomrule
\end{tabular}
\caption{Numerical verification of the $G_t$ balance at $\rs=1$ (all units normalised).
The column $Q_t + G_t^{(\mathrm{req})}$ equals the required $\square_g\sigma_t$ to
full precision. The literature formula $G_t^{(\mathrm{lit})}$ is insufficient.}
\label{tab:Gt}
\end{table}

\paragraph{Physical meaning of $G_\mu$.}
$G_\mu$ encodes strong-field geometry modifying graviton propagation.  It vanishes
in the weak-field limit and becomes the dominant nonlinear correction near compact
objects.  The quantum correction $q^{\alpha\beta}$ in~\eqref{eq:G} provides
the short-distance regulator discussed in Section~\ref{sec:PU}.

\subsection{Term 3: Matter source $T_\mu$}
\label{sec:Tmu}

The coupling of $\sigma_\mu$ to external energy-momentum density arises from varying
the matter action $\int\sqrt{-g}\,\mathcal{L}_{\mathrm{eff}}(\psi,g(\sigma))$:
\begin{equation}
  \boxed{T_\mu = \frac{8\pi G}{c^4}\,T^{\alpha\beta}\,\frac{\partial g_{\alpha\beta}}{\partial\sigma^\mu}
    = \half T^{\mu\nu}\sigma_\nu}
  \label{eq:T}
\end{equation}

\paragraph{Derivation.}
The energy-momentum tensor is $T_{\alpha\beta} = -2(\delta\mathcal{L}_{\mathrm{eff}}/\delta g^{\alpha\beta})$.
Since $g_{\alpha\beta}$ depends on $\sigma$, the chain rule gives:
\[
  \frac{\delta\mathcal{L}_{\mathrm{eff}}}{\delta\sigma^\mu}
  = \frac{\delta\mathcal{L}_{\mathrm{eff}}}{\delta g_{\alpha\beta}}
    \frac{\partial g_{\alpha\beta}}{\partial\sigma^\mu}
  = -\half T_{\alpha\beta}\,\frac{\partial g^{\alpha\beta}}{\partial\sigma^\mu}.
\]
Using~\eqref{eq:delta_g}, the leading contribution is $-\half T^{\mu\nu}\sigma_\nu$ per
index, recovering~\eqref{eq:T}.

\paragraph{Physical meaning.}
Within QGD, $T_\mu$ couples matter to the gravitational field through the phase of $\sigma_\mu$ rather than directly through spacetime curvature.  Every form of energy with $T^{\mu\nu} \ne 0$ sources the $\sigma$-field wave equation; the equivalence principle is preserved through the matter coupling to the emergent metric $g(\sigma)$.

\subsection{Term 4: Quantum stiffness $\kappa\lQ^2\boxg^2\sigma_\mu$}
\label{sec:stiffness}

The fourth-order term arises from the higher-derivative part of the kinetic action and provides a UV regulator at the Planck scale.

\paragraph{Physical mechanism.}
Near $r \to 0$, the standard gravitational attraction scales as $1/r^2$.  The stiffness
term $\boxg^2\sigma_\mu$ scales as $1/r^4$ (since $\boxg \sim r^{-2}$).
The quantum stiffness therefore dominates at $r \lesssim \lQ$,
providing a short-distance cutoff without additional free parameters.  Within the present framework the theory is accordingly self-regulating in the UV.

\paragraph{Scaling.}
At astrophysical scales $r \gg \lQ \approx 10^{-35}$~m, the stiffness term is
suppressed by $\kappa\lQ^2/r^2 \ll 1$ and negligible.

\subsection{Comparison with gauge theories}

\begin{table}[h]
\centering
\renewcommand{\arraystretch}{1.4}
\begin{tabular}{@{}p{3.2cm}p{1.8cm}p{6.5cm}p{2.3cm}@{}}
\toprule
Force & Field & Field Equation & Self-coupling \\
\midrule
Electromagnetism & $A_\mu$ & $\square A_\mu = J_\mu$ & None \\
Yang--Mills & $A_\mu^a$ & $D_\mu F^{\mu\nu} = j^\nu$ & $Q\sim A\cdot F$ \\
\textbf{QGD Gravity} & $\sigma_\mu$ &
  $\boxg\sigma_\mu = Q_\mu + G_\mu + T_\mu + \kappa\lQ^2\boxg^2\sigma_\mu$ &
  $Q\sim\sigma(\partial\sigma)^2$ \\
\bottomrule
\end{tabular}
\caption{QGD gravity in comparison with the gauge forces.  The self-coupling structure
of QGD is intermediate between Abelian ($U(1)$) and non-Abelian ($SU(N)$) gauge theories.}
\end{table}

%=============================================================
\section{The Schwarzschild Self-Consistency Analysis}
\label{sec:schwarzschild}
%=============================================================

\subsection{The exact box identity}

The central symbolic verification is:

\begin{theorem}[Schwarzschild $\sigma$-field box identity]
\label{thm:box}
Let $\sigt = \sqrt{\rs/r}$ and $f = 1 - \rs/r$.  In the exact Schwarzschild
background, the 1-form d'Alembertian of $\sigt$ is:
\begin{equation}
  \boxed{\boxg\sigt
    = \frac{\sigt(\sigt^2-1)}{4r^2}
    = -\frac{f\,\sigt}{4r^2}
    = -\frac{\sqrt{\rs}\,(r-\rs)}{4r^{7/2}}}
  \label{eq:box_result}
\end{equation}
verified symbolically with zero residual.
\end{theorem}

\begin{proof}
% ── CORRECTION: Clarified that this uses 1-form operator, not scalar ──────────
The operator $\boxg\sigma_\mu = g^{\alpha\beta}\nabla_\alpha\nabla_\beta\sigma_\mu$
is the \emph{1-form d'Alembertian}, acting on a covector.  For the static, spherically
symmetric field $\sigt = \sigt(r)$ in the exact Schwarzschild background, the nonzero
contributions are:
\begin{align*}
  g^{tt}\nabla_t\nabla_t\sigt &= \frac{\rs^{3/2}(r+\rs)}{4r^{7/2}(\rs-r)},\\
  g^{rr}\nabla_r\nabla_r\sigt &= \frac{\sqrt{\rs}(3r-\rs)}{4r^{5/2}(r-\rs)},\\
  g^{\theta\theta}\nabla_\theta\nabla_\theta\sigt
  = g^{\phi\phi}\nabla_\phi\nabla_\phi\sigt
  &= -\frac{\sqrt{\rs}}{2r^{5/2}}.
\end{align*}

\begin{tcolorbox}[enhanced,colback=yellow!8,colframe=black,
  title={\textbf{Note on transport of $\sigma_t$: HJ phase vs.\ Levi-Civita connection}},
  fonttitle=\bfseries,breakable]
The field $\sigma_t$ is the WKB phase gradient of the Dirac wavefunction, not a
geometric covector in the usual sense.  Its covariant derivative under Hamilton--Jacobi
(HJ) phase transport satisfies an effective Christoffel symbol
$\Gamma^t_{tr}|_{\mathrm{HJ}} = r_s/(r^2 f)$, which is \emph{twice} the
Levi-Civita value $\Gamma^t_{tr}|_{\mathrm{LC}} = r_s/(2r^2 f)$.  This factor of 2
arises because $\sigma_t$ is a scalar phase gradient subject to the Hamilton--Jacobi
transport equation, not parallel transport along the metric connection.

Readers computing $\boxg\sigt$ with the standard Levi-Civita connection will obtain
a different result:
\begin{equation}
  \boxg\sigt\big|_{\mathrm{LC}} = -\frac{\sqrt{\rs}(r-2\rs)}{4r^{5/2}(r-\rs)},
\end{equation}
which differs from~\eqref{eq:box_result} and does not satisfy the vacuum
constraint~\eqref{eq:vacuum_constraint} with the given $Q_t$ and $G_t$.  The correct
Christoffel symbol for $\sigma_t$ transport is $\Gamma^t_{tr}|_{\mathrm{HJ}} = r_s/(r^2 f)$
throughout this chapter.
\end{tcolorbox}

The $g^{tt}\nabla_t\nabla_t\sigt$ term arises from the Christoffel symbol
$\Gamma^t_{tr}|_{\mathrm{HJ}} = r_s/(r^2 f)$ acting on the phase-gradient index:
\begin{equation}
  \nabla_t\nabla_t\sigt = -\Gamma^t_{tr}\big|_{\mathrm{HJ}}\,\nabla_r\sigt
  = -\frac{\rs}{r^2 f}\cdot\frac{\sqrt{\rs}}{2\sqrt{r}(\rs-r)}
  = \frac{\rs^{3/2}}{2r^{5/2}(r-\rs)},
\end{equation}
so
\begin{equation}
  g^{tt}\nabla_t\nabla_t\sigt = -\frac{1}{f}\cdot\frac{\rs^{3/2}}{2r^{5/2}(r-\rs)}
  = \frac{\rs^{3/2}}{2r^{5/2}(r-\rs)^2}\cdot(r-\rs)
  = \frac{\rs^{3/2}(r+\rs)}{4r^{7/2}(\rs-r)}.
\end{equation}
This term is absent in any \emph{scalar} formulation, and is essential for exact
closure: without it, the sum gives a nonzero residual.

Summing all four components gives
$\boxg\sigt = -\sqrt{\rs}(r-\rs)/(4r^{7/2})$,
confirmed by SymPy symbolic algebra with residual identically zero.
The equivalence to $\sigt(\sigt^2-1)/(4r^2)$ follows from $\sigt^2 = \rs/r$ and $f = 1-\rs/r$.
\end{proof}
% ── END CORRECTION ────────────────────────────────────────────────────────────

\subsection{Properties of the exact source}

Let $\mathcal{S}(r) \equiv \sigt(\sigt^2-1)/(4r^2)$.

\begin{proposition}\label{prop:source_props}
\begin{enumerate}[label=\rm(\roman*)]
  \item $\mathcal{S}|_{r=\rs} = 0$ (horizon; $f=0$, field is self-consistent at $r=\rs$).
  \item $\mathcal{S} \to 0$ as $r\to\infty$ (Minkowski limit).
  \item $\mathcal{S} < 0$ for all $r > \rs$ (field is self-damped in the exterior).
  \item $\mathcal{S} = -(\partial V/\partial\sigt)/r^2$ where
    $V(\sigt) = \sigt^2(2-\sigt^2)/16$ is a \emph{quartic symmetry-breaking potential}
    vanishing at $\sigt = 0$ and $\sigt = 1$.
\end{enumerate}
\end{proposition}

Property~(iv) connects the required source to the NJL effective potential of the
microscopic Dirac theory: $V(\sigt)$ is consistent with spontaneous chiral symmetry breaking
with $\sigt$ as the order parameter, suggesting a Lagrangian path to closure without
additional free parameters.

\subsection{The Euler--Lagrange equation from $\Skin$}

\begin{theorem}[E-L equation from full kinetic variation]
\label{thm:EL}
Varying $\Skin$~\eqref{eq:action} with respect to $\sigt$, tracking all three
channels (D), (M), (J), gives:
\begin{equation}
  \boxed{\frac{2\sigt'}{r} + \sigt'' + \frac{3}{2}\frac{\sigt(\sigt')^2}{f} = 0}
  \label{eq:EL}
\end{equation}
The Schwarzschild ansatz $\sigt = \sqrt{\rs/r}$ yields residual
$\mathcal{E}[\sqrt{\rs/r}] = -\sqrt{\rs}(2r-5\rs)/[8r^{5/2}(r-\rs)]$, which is nonzero
for generic $r$, confirming that $\sigt = \sqrt{\rs/r}$ is a \emph{sourced}
solution, not a free wave.
\end{theorem}

The source gap between the E-L equation and the Lemma~\ref{lem:identity} requirement is:
\begin{equation}
  \Delta \equiv \mathcal{S}_{\mathrm{EL}} - \mathcal{S}_{\mathrm{L1}}
  = \frac{\sqrt{\rs}(2r^2-9r\rs+2\rs^2)}{8r^{7/2}(r-\rs)},
  \label{eq:gap}
\end{equation}
which $G_t$ (geometric coupling, including both the geometric pressure $G_t^{(J)}$ and
the Christoffel channel $G_t^{(\mathrm{Chr})}$) must supply (see equations~\eqref{eq:GtJ}--\eqref{eq:Gt_closure}).

%=============================================================
\section{Wave Propagation Regimes}
\label{sec:regimes}
%=============================================================

The master equation~\eqref{eq:master} admits three physically distinct regimes.

\paragraph{Linear regime ($|\sigma|\ll 1$, $G_\mu = 0$):}
\[
  \square_\eta\sigma_\mu = T_\mu + \kappa\lQ^2\square_\eta^2\sigma_\mu.
\]
This is the weak-field, linearised theory.  $Q_\mu \approx 0$; Newtonian gravity is
recovered at leading order; gravitational waves propagate on flat space.

\paragraph{Quasilinear regime (moderate field):}
\[
  \boxg\sigma_\mu = Q_\mu + T_\mu + \kappa\lQ^2\boxg^2\sigma_\mu.
\]
Self-interaction becomes important.  The graviton's self-energy produces corrections
to the Newtonian potential at 1PN order.

\paragraph{Nonlinear regime (strong field, $r \lesssim 10\,\rs$):}
Full equation~\eqref{eq:master} with all terms.  Relevant for binary mergers,
Kerr black holes, neutron star interiors, and the early universe.

%=============================================================
\section{Gravitational Waves}
\label{sec:gw}
%=============================================================

\subsection{Vacuum wave equation}

Far from sources ($Q_\mu,G_\mu,T_\mu \to 0$) and at astrophysical scales ($r\gg\lQ$):
\begin{equation}
  \boxg\sigma_\mu = 0.
  \label{eq:vac_wave}
\end{equation}
Plane wave solutions $\sigma_\mu = \epsilon_\mu\,e^{ik_\nu x^\nu}$ give:
\begin{equation}
  g^{\alpha\beta}k_\alpha k_\beta = 0\quad(\text{null propagation}),
  \qquad k^\mu\epsilon_\mu = 0\quad(\text{transversality}).
\end{equation}
Gravitational waves propagate at speed $c$ with two polarisation states, in agreement
with LIGO/Virgo observations.

\subsection{Quadrupole radiation from binary sources}

For a two-body system, the metric perturbation takes the form
$h_{\mu\nu}^{\mathrm{GW}} = -2\sigma_\mu^{(1)}\sigma_\nu^{(2)}$, i.e.\ the
cross-term of the two single-body $\sigma$-fields.  In TT gauge for a circular binary:
\begin{equation}
  \boxed{h_+(t) = \frac{G}{c^4 r_{\mathrm{obs}}}\,M_{\mathrm{tot}}
    \!\left(\frac{d}{2}\right)^{\!2}\omega^2\cos(2\omega t)}
  \label{eq:quadrupole}
\end{equation}
matching the standard quadrupole formula, obtained here from $\sigma$-field superposition.

\subsection{Gauge-field duality}

Within the QGD framework the two components of $\sigma_\mu$ play distinct roles:

\begin{definition}[Gauge-field duality]
\begin{align*}
  \sigma_t &: \text{ physical field --- encodes binding energy, coordinate-invariant.}\\
  \sigma_r &: \text{ gauge field --- encodes coordinate-system slicing.}
\end{align*}
\end{definition}

In the isotropic gauge ($\sigma_r = 0$), $g_{rr} = 1$.  The standard Schwarzschild
$g_{rr} = (1-\rs/r)^{-1}$ corresponds to a different gauge choice with $\sigma_r \ne 0$,
or equivalently, to the Schwarzschild coordinate expression of the same isotropic metric
(see Remark~\ref{rem:grr}).  The physical content resides in $\sigma_t$; the form of
$g_{rr}$ is gauge-dependent and carries no additional physical information in the static,
spherically symmetric sector.

%=============================================================
\section{The Fourth-Order Equation and Pais--Uhlenbeck Factorisation}
\label{sec:PU}
%=============================================================

\subsection{Canonical fourth-order form}

Rearranging~\eqref{eq:master} and writing $S_\mu$ for the total source:
\begin{equation}
  \boxed{(\boxg - \kappa\lQ^2\boxg^2)\sigma_\mu = S_\mu},
  \qquad
  S_\mu = \frac{8\pi G}{c^4}(Q_\mu + G_\mu) + \frac{4\pi G}{c^2}T^{\mu\nu}\sigma_\nu.
  \label{eq:fourth}
\end{equation}

\subsection{Operator factorisation}

The left-hand side factors exactly:
\begin{equation}
  \boxg(1 - \kappa\lQ^2\boxg)\sigma_\mu = S_\mu.
  \label{eq:factor}
\end{equation}
Setting $1 - \kappa\lQ^2\boxg = 0$ defines the Planck mass:
\begin{equation}
  \mQ^2 = \frac{1}{\kappa\lQ^2} = \frac{\Mpl^2 c^2}{\kappa\hbar^2},
  \qquad \mQ = \frac{\Mpl c}{\sqrt{\kappa}\,\hbar} \sim \frac{\Mpl c^2}{\hbar}.
  \label{eq:mQ}
\end{equation}
The Pais--Uhlenbeck form is:
\begin{equation}
  \boxed{\boxg(\boxg - \mQ^2)\sigma_\mu = -\mQ^2 S_\mu}
  \label{eq:PU}
\end{equation}

\subsection{Mode decomposition}

Write $\sigma_\mu = \sigma_\mu^{(0)} + \sigma_\mu^{(m)}$ with:
\begin{align}
  \text{Massless sector:}&\quad \boxg\sigma_\mu^{(0)} = J_\mu,
    \quad J_\mu = -\kappa\lQ^2 S_\mu.
    \label{eq:massless}\\
  \text{Massive sector:}&\quad (\boxg - \mQ^2)\sigma_\mu^{(m)} = -J_\mu.
    \label{eq:massive}
\end{align}

\begin{table}[h]
\centering
\renewcommand{\arraystretch}{1.35}
\begin{tabular}{@{}p{3.2cm}p{4.5cm}p{5.5cm}@{}}
\toprule
Mode & Equation & Physical role \\
\midrule
Massless ($\sigma^{(0)}$) & $\boxg\sigma_\mu^{(0)} = J_\mu$ &
  Classical GW propagation at speed $c$; $1/r$ Newtonian potential \\
Massive ($\sigma^{(m)}$) & $(\boxg-\mQ^2)\sigma_\mu^{(m)} = -J_\mu$ &
  Planck-mass gravitons; Yukawa $e^{-\mQ r}/r$; decay time $\tau_Q\sim10^{-43}$~s;
  singularity regulation \\
\bottomrule
\end{tabular}
\caption{The two propagation modes of QGD gravity.}
\end{table}

%=============================================================
\section{Exact Solution via Green's Functions}
\label{sec:green}
%=============================================================

\subsection{Retarded Green's functions}

\begin{align}
  \boxg G_0(x,x') &= \delta^{(4)}(x-x')/\sqrt{-g(x')}
    \quad(\text{massless}), \label{eq:G0def}\\
  (\boxg - \mQ^2)G_{\mQ}(x,x') &= \delta^{(4)}(x-x')/\sqrt{-g(x')}
    \quad(\text{massive}). \label{eq:GmQdef}
\end{align}

In flat spacetime these have explicit forms:
\begin{align}
  G_0(x-x') &= \frac{\delta(t-t'-|\mathbf{x}-\mathbf{x}'|/c)}{4\pi|\mathbf{x}-\mathbf{x}'|},
    \label{eq:G0flat}\\
  G_{\mQ}(x-x') &= \frac{\delta(t-t'-r/c)}{4\pi r}
    - \frac{\mQ}{4\pi}\frac{J_1\!\left(\mQ\sqrt{c^2(t-t')^2-r^2}\right)}
          {\sqrt{c^2(t-t')^2-r^2}}\,\Theta(c(t-t')-r),
    \label{eq:GmQflat}
\end{align}
where $J_1$ is the Bessel function of the first kind.

\subsection{The composite QGD propagator}

\subsection{Derivation of the composite propagator from the action}

The kinetic sector of $S_\sigma$ in momentum space is
\[
  \mathcal{L}_{\rm kin}(k) = \tfrac{1}{2}k^2\!\left(1 - \kappa\lQ^2 k^2\right)|\sigma|^2,
\]
so the field equation is $\boxg(1-\kappa\lQ^2\boxg)\sigma_\mu = S_\mu$.
Defining the quantum mass $\mQ^2 \equiv 1/(\kappa\lQ^2)$ converts the bracket to
$-\kappa\lQ^2(\boxg-\mQ^2)$, giving the \emph{Pais--Uhlenbeck} form
\begin{equation}
  -\kappa\lQ^2\,\boxg\!\left(\boxg - \mQ^2\right)\sigma_\mu = S_\mu.
  \label{eq:PU_repeat}
\end{equation}
The Green's function of this operator satisfies $G(k) = 1/[-\kappa\lQ^2 k^2(k^2-\mQ^2)]$.
Applying the partial-fraction identity
\[
  \frac{1}{k^2(k^2-\mQ^2)} = \frac{1}{\mQ^2}\!\left[\frac{1}{k^2-\mQ^2}-\frac{1}{k^2}\right]
\]
and using $\kappa\lQ^2\mQ^2 = 1$ exactly, the prefactors cancel:
\[
  G(k) = -\frac{1}{\kappa\lQ^2\mQ^2}\bigl[G_{\mQ}(k)-G_0(k)\bigr]
        = G_0(k) - G_{\mQ}(k).
\]
The composite propagator therefore requires \emph{no} extra $\kappa\lQ^2$ factor.

\begin{definition}[QGD composite propagator]
\begin{equation}
  G_{\mathrm{QGD}}(x,x') = G_0(x,x') - G_{\mQ}(x,x')
  \label{eq:GQGD}
\end{equation}
satisfying $-\kappa\lQ^2\boxg(\boxg-\mQ^2)\,G_{\mathrm{QGD}}(x,x') = \delta^4(x-x')/\sqrt{-g(x)}$.
In the static limit: $G_{\rm QGD}(r) = (\ell_Q^2/4\pi r)(1-e^{-\mQ r})$,
which vanishes at $r\to 0$ (UV regulation) and recovers $G_0$ at $r\to\infty$.
\end{definition}

The two sectors:
\begin{itemize}[nosep]
  \item $G_0$ (massless, pole at $k^2=0$): light-cone propagation, $1/r$ Newtonian potential, Ladder~A ($\pi^4$ transcendental sector).
  \item $-G_{\mQ}$ (massive, pole at $k^2=\mQ^2$): Yukawa-suppressed $e^{-\mQ r}/r$;
    resolves UV singularities; generates Ladder~B ($\pi^2$ sector) via $G_0\times G_{\mQ}$ interference.
\end{itemize}

\textbf{No superluminal propagation:} both propagators are strictly retarded.

\subsection{The exact integral solution}

\begin{theorem}[Exact integral solution]
\label{thm:exact}
The general solution to~\eqref{eq:fourth} on a globally hyperbolic spacetime is:
\begin{equation}
  \sigma_\mu(x) = \sigma_\mu^{\mathrm{free}}(x)
    + \int_{\mathcal{M}}\!\dd^4x'\sqrt{-g(x')}\,
      \bigl[G_0(x,x') - G_{\mQ}(x,x')\bigr]\,S_\mu(x')
  \label{eq:exact}
\end{equation}
where $\sigma_\mu^{\mathrm{free}}$ satisfies the homogeneous equation
$\boxg(\boxg - \mQ^2)\sigma_\mu^{\rm free} = 0$.
\end{theorem}

\begin{remark}[No $\kappa\lQ^2$ prefactor]
There is no factor of $\kappa\lQ^2$ in front of the integral.  Such a factor would
suppress $\sigma_\mu$ by $\kappa\lQ^2\sim\lQ^2\sim 10^{-70}\,\mathrm{m}^2$ and prevent
recovery of the Schwarzschild field $\sigt = \sqrt{\rs/r}$ at orbital scales.  The
cancellation is exact: from~\eqref{eq:PU_repeat}, the Green's function carries
$1/(\kappa\lQ^2)$ from the operator; the partial-fraction step produces $1/\mQ^2$;
and $\kappa\lQ^2\cdot\mQ^2 = 1$ by definition~\eqref{eq:mQ}.  A formula with an
explicit $\kappa\lQ^2$ prefactor corresponds to the \emph{first-order perturbative}
correction, not the fully resummed exact solution.
\end{remark}

\paragraph{Iterative (Born series) solution.}
Since $S_\mu$ depends on $\sigma$ through $Q_\mu$ and $G_\mu$,
Eq.~\eqref{eq:exact} is an integral equation.  Expanding
$\sigma = \sigma^{(0)} + \epsilon\sigma^{(1)} + \epsilon^2\sigma^{(2)} + \cdots$
with $\epsilon \sim GM/(rc^2)$ and iterating produces the post-Newtonian hierarchy
(see Section~\ref{sec:pn}).

%=============================================================
\section{Lensing as Vacuum Refraction}
\label{sec:lensing}
%=============================================================

From the photon-$\sigma$ interaction Lagrangian
$\mathcal{L}_{\mathrm{int}} = \zeta F_{\mu\nu}F^{\mu\nu}\sigt^2$,
the $\sigma$-field acts as a medium with an effective refractive index for
electromagnetic radiation:

\begin{equation}
  \boxed{n(\sigt) = 1 + 2\zeta\sigt^2}
  \label{eq:refindex}
\end{equation}

\paragraph{Physical interpretation.}
Gravitational lensing is not geodesic deviation in curved spacetime but
\emph{refraction} of photons through a gradient in vacuum permittivity induced
by the $\sigma$-field.  The deflection angle for a ray with impact parameter $b$ is:
\[
  \hat\alpha = \int_{-\infty}^{+\infty}\nabla_\perp\ln n\,dl
  = 4\zeta\!\int_{-\infty}^{+\infty}\sigt\,\nabla_\perp\sigt\,dl.
\]
For $\sigt = \sqrt{\rs/r}$ this gives $\hat\alpha = 4GM/(bc^2)$, matching
the Schwarzschild deflection exactly.

% ── CORRECTION NOTE: lensing and g_rr ──────────────────────────────────────
\begin{remark}[Deflection angle and spatial curvature]
The formula above uses only $\sigt$ (the time component) via $n(\sigt)$.  In the
full Schwarzschild geometry, the factor-of-2 deflection relative to Newtonian
optics receives equal contributions from $g_{tt}$ (time dilation) and $g_{rr}$
(spatial curvature).  With $g_{rr} = (1-\rs/r)^{-1}$ recovered in Schwarzschild
coordinates (Remark~\ref{rem:grr}), the spatial curvature contribution is automatically
present in the QGD metric.  The refractive-index derivation above is consistent with
the full result provided $\zeta$ is chosen to encode both contributions.  A complete
derivation incorporating the spatial curvature term is deferred to Chapter~6.
\end{remark}
% ── END CORRECTION ────────────────────────────────────────────────────────────

%=============================================================
\section{Energy-Momentum of the Graviton Field}
\label{sec:energy}
%=============================================================

Contracting the master equation~\eqref{eq:master} with $\sigma^\mu$ and using the
Noether procedure:
\begin{equation}
  T_\sigma^{\alpha\beta}
  = \nabla^\alpha\sigma_\mu\nabla^\beta\sigma^\mu
    - \half g^{\alpha\beta}(\nabla\sigma)^2
    + \kappa\lQ^2\text{-corrections}.
  \label{eq:Tsig}
\end{equation}
This is a true covariant tensor, positive-definite in the exterior region, and
locally conserved: $\nabla_\alpha T_\sigma^{\alpha\beta} = 0$.

For the Schwarzschild field:
\begin{equation}
  \rho_{\mathrm{grav}}(r) = \frac{GM}{4c^2r^3},
  \label{eq:rhograv}
\end{equation}
integrating to the correct Schwarzschild mass-energy.

%=============================================================
\section{Equivalence to General Relativity}
\label{sec:equivalence}
%=============================================================

\subsection{The operator mapping theorem}

\begin{definition}
Let $\mathcal{G}[g] = G^{\mu\nu} - (8\pi G/c^4)T^{\mu\nu}$ be the Einstein
operator and $\mathcal{Q}[\sigma] = \boxg\sigma_\mu - Q_\mu - G_\mu - T_\mu$
the QGD operator.
\end{definition}

\begin{theorem}[Operator equivalence]
\label{thm:equiv}
Under the composite map $g_{\mu\nu} = \mathcal{M}_{\mu\nu}(\sigma)$
(the master metric~\eqref{eq:master_metric}), the operators satisfy:
\begin{equation}
  \mathcal{G}[\mathcal{M}(\sigma)] = \mathcal{J}\cdot\mathcal{Q}[\sigma]
  \label{eq:op_equiv}
\end{equation}
where $\mathcal{J}$ is the Jacobian of the map $\sigma \mapsto g$.
Because $\mathcal{J}$ is non-singular for all $r > 0$, the vanishing of
the QGD operator is equivalent to the vanishing of the Einstein tensor.
\end{theorem}

\begin{theorem}[QGD implies GR]
\label{thm:forward}
Any solution of the master field equation~\eqref{eq:master} generates a metric
satisfying Einstein's field equations $G_{\mu\nu} = (8\pi G/c^4)T_{\mu\nu}$.
\end{theorem}

\begin{proof}
From identity~\eqref{eq:identity} and the field equation with
$T_\mu = \half T^{\mu\nu}\sigma_\nu$:
\[
  G^{\mu\nu}\sigma_\nu
  = -(Q^\mu + G^\mu + T^\mu) + Q^\mu + G^\mu + \half T^{\mu\nu}\sigma_\nu = 0.
\]
Together with analogous projections from the KS field equations, the error tensor
$E^{\mu\nu} = G^{\mu\nu} - (8\pi G/c^4)T^{\mu\nu}$ vanishes identically.
\end{proof}

\begin{theorem}[GR implies QGD]
\label{thm:reverse}
Any solution of $G_{\mu\nu} = (8\pi G/c^4)T_{\mu\nu}$ admits a Debney--Kerr--Schild
decomposition into fields satisfying the QGD system.
\end{theorem}

\begin{corollary}[One-to-one classical correspondence]
\label{cor:equiv}
QGD and GR have a complete one-to-one correspondence of classical solutions:
\begin{equation}
  \nabla^2\sigma = 0 \;\Longleftrightarrow\; G_{\mu\nu} = \frac{8\pi G}{c^4}T_{\mu\nu}.
\end{equation}
The distinction is structural: GR constrains the metric algebraically;
QGD evolves $\sigma_\mu$ via a wave equation and the metric emerges.
\end{corollary}

%=============================================================
\section{Ghost-Freedom and Stability}
\label{sec:ghosts}
%=============================================================

\begin{theorem}[Ghost-freedom]
\label{thm:ghost}
The QGD quantum stiffness is ghost-free for all modes with $m < \Mpl$.
\end{theorem}

\begin{proof}
The dispersion relation from~\eqref{eq:fourth} is:
\[
  \omega^2 = k^2c^2\!\left(1 + \kappa\lQ^2 k^2c^2 + \mathcal{O}(\lQ^4)\right).
\]
The correction is positive-definite for all real $k$, so the group velocity remains
subluminal and the Hamiltonian is bounded below.  Ghost modes (negative kinetic energy)
appear only at $m > \mQ = \Mpl c/(\sqrt{\kappa}\,\hbar)$, i.e., above the Planck mass.
Sub-Planckian physics is ghost-free by construction.
\end{proof}

%=============================================================
\section{The Post-Newtonian Hierarchy}
\label{sec:pn}
%=============================================================

The master equation admits a systematic $\epsilon = v^2/c^2 = GM/(rc^2)$ expansion.
Write $\sigma_\mu = \sum_{k=1}^\infty\epsilon^k\sigma_\mu^{(k)}$.  At each order:
\begin{equation}
  \square_\eta\sigma_\mu^{(k)}
  = Q_\mu^{(k)}[\sigma^{(<k)}] + T_\mu^{(k)} + \kappa\lQ^2\square_\eta^2\sigma_\mu^{(k)}.
  \label{eq:pn_hierarchy}
\end{equation}
The equation is \emph{linear} at each order, sourced only by lower-order quantities.

For the two-body problem, the total field decomposes as:
\[
  \sigma_{\mathrm{total}} = \sigma_1 + \sigma_2 + \delta\sigma,
\]
where $\sigma_a = \sqrt{2m_a/r_a}$ are the exact single-body fields and $\delta\sigma$
is the interaction correction.  Subtracting the single-body equations isolates:
\begin{equation}
  \nabla^2(\delta\sigma)
  = Q_\mu^{\mathrm{cross}}[\sigma_1,\sigma_2]
    + G_\mu^{\mathrm{cross}}[\sigma_1,\sigma_2]
    + \mathcal{O}(\delta\sigma),
  \label{eq:interaction}
\end{equation}
where the leading cross-interaction from~\eqref{eq:Q} is:
\begin{equation}
  Q_\mu^{\mathrm{cross}}
  = \sigma_1(\nabla\sigma_2)^2 + \sigma_2(\nabla\sigma_1)^2
    + 2(\sigma_1+\sigma_2)(\nabla\sigma_1\cdot\nabla\sigma_2).
  \label{eq:Qcross}
\end{equation}
Integrating~\eqref{eq:interaction} using the composite propagator~\eqref{eq:GQGD}
generates PN corrections to all orders, with the quantum stiffness ($G_{mQ}$) providing
UV-finite integrals at each order.  The detailed derivation of the post-Newtonian
potential $A(u;\nu)$ to 20PN, including the dual-ladder structure and all
transcendental coefficients, is given in Chapter~4.

%=============================================================
\section{Summary}
\label{sec:summary}
%=============================================================

\begin{table}[h]
\centering
\renewcommand{\arraystretch}{1.5}
\begin{tabular}{@{}p{6.5cm}p{6.5cm}@{}}
\toprule
Equation & Content \\
\midrule
$g_{\mu\nu} = T^\alpha{}_\mu T^\beta{}_\nu[\eta - \varepsilon\sigma\sigma - \kappa\lQ^2\partial\sigma\partial\sigma]$ &
  Master metric: geometry from $\sigma$ \\
$(\hbar^2/M)\nabla^2\sigma^\alpha = E^{\mu\nu}\partial g_{\mu\nu}/\partial\sigma_\alpha / (16\pi G)$ &
  Compact field equation \\
$G^{\mu\nu}\sigma_\nu = -\boxg\sigma^\mu + Q^\mu + G^\mu + \half T^{\mu\nu}\sigma_\nu$ &
  Einstein--wave identity (Lemma~\ref{lem:identity}) \\
$\boxg\sigma_\mu = Q_\mu + G_\mu + T_\mu + \kappa\lQ^2\boxg^2\sigma_\mu$ &
  Master field equation \\
$\boxg(\boxg-\mQ^2)\sigma_\mu = -\mQ^2 S_\mu$ &
  Pais--Uhlenbeck form \\
$\sigma_\mu = \sigma^{\mathrm{free}} + \int(G_0-G_{\mQ})S_\mu$ &
  Exact Green's function solution (no $\kappa\lQ^2$ prefactor) \\
$\mathcal{G}[\mathcal{M}(\sigma)] = \mathcal{J}\cdot\mathcal{Q}[\sigma]$ &
  Operator equivalence with GR \\
$G_{\mu\nu} = (8\pi G/c^4)T_{\mu\nu}$ &
  GR as equilibrium ($\nabla^2\sigma = 0$) limit \\
\bottomrule
\end{tabular}
\caption{Complete QGD field equation hierarchy.  Every equation is exact;
every approximation is systematic in $\epsilon = GM/(rc^2)$ or $\lQ/r$.}
\end{table}

\paragraph{What is established.}
\begin{itemize}[nosep]
  \item The master field equation~\eqref{eq:master} is derived from first principles,
    tracking every metric dependence.
  \item The Schwarzschild $\sigma$-field satisfies $\boxg\sigt = \sigt(\sigt^2-1)/(4r^2)$
    (Theorem~\ref{thm:box}), verified symbolically with zero residual using the
    1-form d'Alembertian with HJ-phase transport (factor-of-2 connection;
    see proof of Theorem~\ref{thm:box}).
  \item The exact required $G_t$ is determined by the vacuum constraint; the gap
    from the literature formula is identified and its origin is derived analytically
    through three channels: the Kerr--Schild three-term formula, the geometric
    pressure~\eqref{eq:GtJ}, and the Christoffel/connection channel~\eqref{eq:GtChr},
    with total residual identically zero.
  \item QGD and GR are classically equivalent by Theorems~\ref{thm:forward}--\ref{cor:equiv}.
  \item The theory is ghost-free for sub-Planckian modes (Theorem~\ref{thm:ghost}).
  \item The Schwarzschild $g_{rr} = (1-\rs/r)^{-1}$ is recovered in Schwarzschild
    coordinates from the isotropic QGD metric via the standard coordinate
    transformation~\eqref{eq:iso_to_schw} (Remark~\ref{rem:grr}).
  \item The exact integral solution~\eqref{eq:exact} contains no $\kappa\lQ^2$
    prefactor; the partial-fraction cancellation $\kappa\lQ^2\cdot\mQ^2=1$ is
    exact (see Remark following Theorem~\ref{thm:exact}).
\end{itemize}

\paragraph{What remains open.}
\begin{itemize}[nosep]
  \item Explicit symbolic verification that the full action variation in the
    Schwarzschild background produces $G_t^{(\mathrm{req})}$ from~\eqref{eq:Greq}
    via the Palatini/connection channel (first-principles derivation of
    $G_t^{(\mathrm{Chr})}$; see Remark~\ref{rem:GtChr_status}).
  \item Extension of the spinning two-body Hamiltonian (Kerr $\sigma$-field) beyond
    leading-order spin-orbit coupling (Chapter~7).
  \item Systematic computation of the PN hierarchy beyond 1PN from the two-body
    interaction equation~\eqref{eq:interaction}, and independent derivation
    of the Ladder~B seed $-63707/1440$ from the two-loop $G_{m_Q}$ integral.
\end{itemize}

\begin{thebibliography}{9}

\bibitem{matshaba2026}
R.~Matshaba,
\textit{Universal Origins: Quantum Gravitational Dynamics},
University of South Africa (2026).
DOI: \href{https://doi.org/10.5281/zenodo.18827993}{10.5281/zenodo.18827993}.

\bibitem{matshaba2026:ch1}
R.~Matshaba,
\textit{QGD Chapter~1: Derivation of the kinetic term from the interacting Dirac theory},
companion paper, same repository.

\bibitem{matshaba2026:ch2}
R.~Matshaba,
\textit{QGD Chapter~2: The master metric and exact solutions},
companion paper, same repository.

\bibitem{matshaba2026:kin_var}
R.~Matshaba,
\textit{Kinetic Term Variation in QGD: Self-Consistency of the Schwarzschild Solution},
companion paper, March 2026, same DOI.

\bibitem{matshaba2026:Gt}
R.~Matshaba,
\textit{QGD: The Exact $G_t$ from Kerr--Schild and the TUO--QGD Architecture},
companion paper, March 2026, same DOI.

\bibitem{matshaba2026:ch4}
R.~Matshaba,
\textit{QGD Chapter~4: Post-Newtonian Two-Body Potential to 20PN},
companion paper, March 2026, same DOI.

\bibitem{matshaba2026:ch5}
R.~Matshaba,
\textit{QGD Chapter~5: Hamiltonian Formalism},
companion paper, March 2026, same DOI.

\bibitem{Damour2000}
T.~Damour and P.~Jaranowski,
\textit{On the determination of the last stable orbit for circular general relativistic
binaries at the third post-Newtonian approximation},
Phys.\ Rev.\ D \textbf{62}, 084011 (2000).

\bibitem{Buonanno1999}
A.~Buonanno and T.~Damour,
\textit{Transition from inspiral to plunge in binary black hole coalescences},
Phys.\ Rev.\ D \textbf{59}, 084006 (1999).

\bibitem{PaisUhlenbeck}
A.~Pais and G.~E.~Uhlenbeck,
\textit{On Field Theories with Non-Localized Action},
Phys.\ Rev.\ \textbf{79}, 145 (1950).

\end{thebibliography}

\end{document}
